\section{Praktikum Orientasi}
\label{sec:praktikum-bab1}

\begin{praktikumbox}
\textbf{Judul:} Penyiapan portofolio, jurnal praktikum, dan target
belajar pribadi.\\[0.3em]
\textbf{Sub-CPMK:} persiapan untuk 1.1--4.3 \quad\textbullet\quad
\textbf{Alokasi waktu:} 50 menit \quad\textbullet\quad
\textbf{Bentuk:} individu
\end{praktikumbox}

\subsection{Tujuan Praktikum}

Menyiapkan seluruh perangkat administratif dan teknis yang
menunjang pembelajaran sepanjang semester, serta menyusun target
belajar pribadi yang terukur.

\subsection{Alat dan Bahan}

\begin{itemize}
    \item Buku ajar ini beserta dokumen RPS mata kuliah.
    \item Komputer dengan TASM, TLINK, DEBUG, emu8086, DOSBox, dan
          Turbo Debugger.
    \item Direktori kerja, misalnya \file{D:\textbackslash asm}.
\end{itemize}

\subsection{Langkah Kerja}

\begin{enumerate}
    \item \textbf{Menyiapkan struktur direktori kerja.}
    \begin{enumerate}
        \item Buat direktori \file{D:\textbackslash asm}.
        \item Di dalamnya buat subdirektori
              \file{bab03}, \file{bab04}, \ldots{} \file{bab16}.
        \item Setiap program yang Anda tulis disimpan pada
              subdirektori bab yang bersesuaian. Portofolio inilah
              yang dikumpulkan pada akhir semester.
    \end{enumerate}

    \item \textbf{Menyiapkan jurnal praktikum.}
    \begin{enumerate}
        \item Buat berkas \file{jurnal.md} atau buku catatan
              terpisah.
        \item Salin format pada Tabel~\ref{tab:format-jurnal} menjadi
              kerangka jurnal.
        \item Isi satu baris jurnal untuk sesi hari ini sebagai
              latihan.
    \end{enumerate}

    \item \textbf{Memeriksa perangkat praktik.}
    \begin{enumerate}
        \item Periksa ketersediaan TASM, TLINK, dan DEBUG di dalam
              DOSBox.
        \item Periksa ketersediaan emu8086 dan Turbo Debugger.
        \item Baca Bab~VII untuk alur perakitan, pengaitan, dan
              penjalanan program.
    \end{enumerate}

    \item \textbf{Menelusuri buku ajar.}
    \begin{enumerate}
        \item Buka Daftar Isi dan temukan bab yang bersesuaian dengan
              minggu ke-1, ke-8, dan ke-16.
        \item Temukan Lampiran~C (glosarium) dan Lampiran~A (rubrik
              penilaian). Catat halaman masing-masing.
        \item Baca rubrik penilaian praktikum pada Lampiran~A, lalu
              tuliskan satu hal yang akan Anda lakukan agar mencapai
              tingkat tertinggi pada rubrik tersebut.
    \end{enumerate}

    \item \textbf{Menyusun target belajar pribadi.}
    \begin{enumerate}
        \item Tuliskan tiga target belajar untuk semester ini dengan
              rumusan yang terukur. Contoh rumusan yang benar:
              ``Saya dapat menulis program pengurutan assembly
              tanpa melihat contoh pada akhir minggu ke-11.''
        \item Tuliskan satu kebiasaan belajar yang akan Anda jaga,
              misalnya ``menjalankan setiap contoh program dalam
              waktu 24 jam setelah kelas.''
    \end{enumerate}
\end{enumerate}

\subsection{Hasil yang Dikumpulkan}

\begin{enumerate}
    \item Struktur direktori kerja (bukti tangkapan layar).
    \item Berkas jurnal praktikum dengan satu entri terisi.
    \item Lembar target belajar pribadi.
    \item Catatan halaman Daftar Isi, glosarium, dan rubrik penilaian.
\end{enumerate}

\begin{asesmenbox}
\textbf{Asesmen orientasi (bukan komponen nilai, bersifat
prasyarat).} Praktikum ini tidak dinilai secara numerik, tetapi
menjadi \emph{prasyarat} untuk mengikuti praktikum
Bab~\ref{chap:pengantar-assembly}.
Mahasiswa yang belum memiliki struktur direktori dan jurnal tidak
dapat mengumpulkan Tugas 1.

\begin{tabular}{@{}p{7.0cm}p{6.4cm}@{}}
    \toprule
    \textbf{Butir} & \textbf{Kriteria terpenuhi} \\
    \midrule
    Struktur direktori & Subdirektori bab dibuat lengkap dan
    penamaannya konsisten \\
    \addlinespace
    Jurnal praktikum & Format sesuai tabel dan minimal satu entri
    terisi \\
    \addlinespace
    Target belajar & Tiga target terukur dan satu kebiasaan belajar
    dituliskan \\
    \bottomrule
\end{tabular}
\end{asesmenbox}
